\section{First-missing-seven ORIGINAL coefficients by a finite source patch}
\begin{theorem}
Every final-deletion-closed class of ACTUAL normalized increasing
arms with FULL C and first missing weight at most SEVEN has FULL
pair A2, and ORIGINAL UNAMPLIFIED coefficient unimodality of the
strict core and BOTH relatives at EVERY strict central peak.
ONE arbitrary legal heavier append is allowed.
For two nonrectangular members, every source window of length
at least FOUR has unimodal positive part with initial positive support.
Raw-source UI and three-window UI are NOT asserted; both can fail.
\end{theorem}
The general conjecture remains OPEN. This is a genuine coefficient
bridge for a larger first-missing range, not a cumulative-LC substitute.

\subsection{The complete mixed-owner reduction}
The frozen first-missing-six criterion handles the deletion-closed
subclass with $\kappa\le6$, including its arbitrary heavier append.
Consider an old pair with maximum first missing weight seven,
assuming FULL OLD A2 in the count induction. A1 and the both-rising
FULL-C and both-nonrising actual resources are unchanged. Only the
mixed source Q comparison can fail.

If its rising opposite has $\kappa\le6$, reuse the frozen pointwise
classification. A $\kappa\le6$ one-part exception would require
the other arm to be $(1,2)$, so it cannot occur in a pair with
maximum first missing weight seven. The new opposite therefore
has $\kappa=7$. Write its weight $b\ge21$.

For a nonincreasing owner with at least FOUR non-one parts, the
OLD sharp gap-six theorem applies since $\alpha\ge1=7-6$.
For an owner with two or three parts use the exact initial interval
exchange to take $\alpha_* =\max(\alpha,t_1-1)$.
It preserves its weight and arm polynomial, and preserves sortedness
because the exchanged first non-one part is at most the original one.
If $\alpha_*\ge2$, the OLD gap-five theorem applies.
The remaining canonical owner has $\alpha=1$, $t_1=2$, and
two or three non-one parts.

Under a putative positive mixed recovery, the OLD arbitrary-gap
finite-owner theorem gives, with $d=\kappa-\alpha=6$,
\[
                     D=\sum(t_i-1)\le2d(d-3)=36.       \tag{1}
\]
Hence EVERY such owner is exactly one of the 341 sorted words
\[
 (2,T),\ 2\le T\le36;\qquad
 (2,L,T),\ 2\le L\le T,\ L+T\le37.
\]
All negative reverse indices are already paid by the first birth:
at $r=-1$, their weight is at least five, so $a+2-7\ge0$.

At $r\ge0$, let $s_r=\Delta q_Z(r)$ and
$C=s_{r+1}-s_r>0$ with $s_r\ge0$.
The OLD early source theorem gives the TWO compulsory payments
\[
 \Delta(E_YF_Z)_{j+1}\ge F_Z(a-r+1-7)+F_Z(a-r+1-9),       \tag{2}
\]
because the first source increment is at least one and its increment
at first weight plus two is at least one. The source prefix is
nondecreasing in the mixed branch; every other term is nonnegative.
The complete certificate reconstructs EVERY owner and EVERY such
curvature edge. There are 853 edges, 851 paid by (2), and exactly TWO
unpaid owner/rank combinations:
\[
       Z=(1,2,2),\ r=0,\ C=1;\qquad
       Z=(1,2,2,2),\ r=0,\ C=2.                    \tag{3}
\]
ALL 71 distinct nonnegative payment margins are checked in one
ordinary kernel aggregate. This is a complete finite complement
of (1), not sampled opposite arms or an unbounded numerical premise.

For a ONE-part owner the positive primitive is an indicator, so only
$r=-1$ can have positive curvature. Its birth is missing only at
weight $a=3$ or four. The weight-three owner is $(1,2)$; weight
four has the equivalent forms $(1,3)$ and $(1,1,2)$.
Thus (3) and these two one-part polynomials exhaust ALL new
post-LAST-mode Q exceptions.

\subsection{Every exceptional source, with its full protected prefix}
Put
\[
 \delta=(E_Y)_7\ge1,\quad\eta=(E_Y)_8\ge\delta,
 \quad\zeta=(E_Y)_9\ge\eta+1.
\]
The last inequality is the OLD mode-free early source curvature,
not a differentiated coefficient bound. Ranks through six of $F_Y$
are full, $(F_Y)_7=y\in[1,63]$, and $\delta=64-y$.
For each exception the exact source $E=(2-z)q_Xq_Y$ has the
following form. Its prefix $c=(E_0,\ldots,E_L)$ is positive
integer and UI; $E_0=2$.
\begin{center}
\begin{tabular}{c|c|c|c|c}
owner & $L$ & entries after $L$ through $R$ & $R$ & protected body\\\hline
$(1,2)$ & $b-3$ & $0,0,1,1$ & $b+1$ & $E_j\ge1$, $1\le j\le b-3$\\
one-part weight four & $b-2$ & $0,1,1,0$ & $b+2$ & $E_j\ge1$, $1\le j\le b-2$\\
$(1,2,2)$ & $b-1$ & $1,2,0$ & $b+2$ & $E_j\ge2$, $1\le j\le b-2$\\
$(1,2,2,2)$ & $b-1$ & $3,4,2$ & $b+2$ & $E_j\ge4$, $1\le j\le b-1$
\end{tabular}
\end{center}
Every source entry after the displayed $R$ is nonpositive.
The last row is needed ONLY when $\eta=\delta=1$; otherwise
its only possible recovery has size $2-\eta\le0$ and Q holds.
The integer implication $1\le\delta\le\eta<2\Rightarrow\delta=\eta=1$
is a new kernel.

Here are the full mode and lower-body justifications, rather than
assuming that a short tested pattern controls the whole prefix.
For the first three rows, the first source value at $L$ is $\delta$.
The forward product coefficients at ranks six and seven are,
respectively,
\[
 (68,y+72),\quad(74,y+84),\quad(83,y+102).
\]
All have a strict forward increase, so the reverse LAST ordinary
mode is before $L$. In row three, $E_{L-1}=\delta+\eta\ge2$;
UI of the prefix and $E_0=2$ give its stated body bound even
though its LAST entry $E_L=\delta$ may be only one.

In row four, $\delta=\eta=1$ implies
$(F_Y)_7=63$, $(F_Y)_8=125$. The product has
$P_8=373$, $P_9=(F_Y)_9+493>373$ since $b\ge21$ puts rank
nine in the actual support. Thus its LAST ordinary mode is
before $L=b-1$. Here $E_L=\zeta+3\ge5$.
Also $q_X(1\hbox{-coefficient})=4$ and $q_Y(1\hbox{-coefficient})\ge6$
because $\kappa_Y\le h_Y+2$. Hence $E_1\ge19$.
UI of the positive prefix gives $E_j\ge4$ on $1\le j\le L$;
the initial coefficient $E_0=2$ is NOT incorrectly included.

In ALL rows, FULL OLD A2 supplies source-prefix UI through the
ordinary first mode; the whole plateau joins by the OLD argument.
Q holds thereafter except at the classified later edge. The positive
value at $L$ and A1 signs preclude an earlier negative or zero
in this tail. Thus the ENTIRE prefix $c$ has the shape used above.
The upper tails are obtained directly from the corresponding exact
owner signed primitive, since all later source reads are before
the opposite first source. No polynomial source is discarded.

\subsection{A general finite-height and finite-width source repair}
Let $E$ be finite and have a positive UI prefix on $[0,L]$ and be nonnegative
through $R\ge L$, then nonpositive after $R$. Suppose for some
$B,H$ that
\[
 E_j\ge H\quad(1\le j\le B),\qquad
 E_n\le H\quad(L<n\le R).
\]
For a window length $t\ge1$ satisfying $R\le t+B$, its positive
uniform window $([t]E)^+$ is globally UI with initial positive support.

Through $L$ it agrees with a uniform convolution of the UI prefix.
Before the first outgoing edge, $n<t$, every difference is
the nonnegative incoming coefficient. For critical $L<n\le R$
with $n>t$, its outgoing index satisfies $1\le n-t\le B$,
so $E_n-E_{n-t}\le0$.
The ONE edge $n=t$ is deliberately unconstrained: if it falls,
all subsequent critical edges are paid; if it rises, there was
no earlier fall because all previous edges were before $t$.
After $R$ the signed-window tail is nonincreasing in positive
part until its whole input lies in the nonpositive tail, then zero.
This proves a global no-fall-then-rise statement. Finite support
and a positive initial value give a whole up/down mode and no holes.

GeneralFiniteSourcePatch.lean verifies the arbitrary-width splice,
the ENTIRE body-guarded critical payment and the whole finite
positive-window up/down theorem, including clipping and the first
outgoing edge. It reuses the frozen signed-tail and finite-mode
kernels of KappaSixWindowRepair. To rebuild this import, compile
that frozen module to an olean first; its frozen source/log is unchanged.
The actual interval convolution and head-UI identification are
the written assembly above, not a cumulative-LC implication.

For the four rows take, respectively,
\[
 (L,R,B,H)=(b-3,b+1,b-3,1),\ (b-2,b+2,b-2,1),
 \ (b-1,b+2,b-2,2),\ (b-1,b+2,b-1,4).
\]
EVERY $t\ge4$ satisfies $R\le t+B$. The complete Nat read guards
are kernel checked. The first-outgoing exception is essential in
the fourth row because its initial coefficient is two, below four.

The restriction four cannot be replaced by three in this universal
window assertion: the ACTUAL FULL-C pair $(1,2)$ and $(1,6^5)$,
whose nonbinary weight is 30, has three consecutive coefficients
of $[3]E$ equal to $(3,1,2)$. This is a source-window obstruction,
not a counterexample to A2 or ORIGINAL coefficient unimodality.

\subsection{Legal heavier appends, full A2 and original polynomials}
At every new exceptional pair, the small owner has weight at most
seven and the opposite weight is $b\ge21$. Appending to the small
owner and making it heavier forces $t\ge b-a\ge14$.
Appending to the first-missing-seven arm forces sorted legality
$t\ge\max Y\ge6$. Thus EVERY required new exceptional window
has length at least FOUR. The whole two-case guard is a new kernel.

Now repeat the FULL A2 count induction of the preceding criterion.
Old pairs with maximum first missing weight at most six use that
frozen theorem and its arbitrary-heavy-append corollary. At an old
pair with maximum seven, normal Q gives the OLD whole source cone;
the four classified exceptions use the repaired window above.
The heavier NEW first ordinary mode is before its new arm weight,
so its whole rising source prefix equals this legal window.
The weighted-difference sign identity and terminal zeros give
NEW full A2. Rectangles and heavier one-part bases are OLD.
No unrestricted raw-source UI is assumed.

After this induction, every four-or-longer source window has the
stated whole shape. The same legal-window argument permits ONE
arbitrary heavier append from the class without assuming C or a
small first missing weight of the appended arm. The OLD least-peak,
proper-relative, and reflection joins, within the final-deletion-closed
family, prove ORIGINAL strict-core and BOTH-relative coefficient UI
at EVERY strict central peak. No amplification is used.

\subsection{An explicit two-dominant-jump original family}
Adjoin to $\mathcal A_{12}$:
\begin{enumerate}
\item ALL $1^\alpha2^u3^r4^sT_1T_2$ with $W_0=3r+4s$,
      $T_1\ge\max(5,2W_0)$ and $T_2\ge2(W_0+T_1)$;
\item ALL $1^\alpha2^uV$ with capacities in $\{3,4,5,6\}$
      and total nonbinary weight at most 32.
\end{enumerate}
Call the union $\mathcal A_{13}$. Complete capacity-four C and
the OLD iterated dominant-append theorem give FULL C uniformly
in every binary count and initial padding in item one. Item two
uses the OLD complete weight-32 theorem. Each last append raises
first missing weight by at most one, so every member has $\kappa\le7$.
Deleting a last part stays in the union, including all bases.
The criterion therefore proves FULL A2 and ORIGINAL core/BOTH
relative coefficient UI at every strict peak. BOTH arms may have
TWO independent dominant jumps with no weight-order restriction;
ONE further arbitrary legal heavier append is also allowed.

\paragraph{Exact verification scope.}
Three new generic finite-patch reports, five new complete source/owner/
mode/read/legal-window reports, and ONE ordinary kernel aggregate of
all 71 distinct finite owner margins make NINE new reports.
The finite premise is explicitly the entire 341-owner domain,
853 curvature edges, 851 paid and TWO separately handled exceptions.
Independent suffix boxes reconstruct EVERY state, curvature edge,
exception and aggregate obligation. Separate actual source, window,
full-A2 and true greedy original-ideal computations are QA only.
Old initial exchange, gap-five/gap-six and their already declared
finite complements, arbitrary-gap finite-owner bounds, source
curvature, actual C/A1/Q, LC-independent uniform UI convolution,
and whole A2/original generation are reused or written above.
This is not an end-to-end Lean composition formalization.
General higher first missing weights, higher-capacity C and the
GENERAL original unamplified conjecture remain OPEN.
